\section{Introduction}
\label{sec:introduction}

Long-wavelength limits of integrable spin chains relate quantum lattice
models to classical field theories.  Coherent states associate a smooth
spin field with a slowly varying lattice configuration, while the
quantum inverse-scattering construction supplies the lattice Lax
operators.  Systematic continuum constructions recover the
Landau--Lifshitz model and its isotropic and anisotropic generalizations
\cite{AvanDoikouSfetsos2010,DoikouKaraiskos2011}.

We consider the Class 5 and Class 6 spin-$\tfrac12$ chains, two
non-Hermitian deformations of the Heisenberg XXX chain in the
classification of Ref.~\cite{deLeeuwPribytokRyan2019}.
Their deformed Landau--Lifshitz continuum models were obtained in
Ref.~\cite{deLeeuwFontanellaNietoGarcia2026}, where a boost construction
generates higher conserved charges and the construction of Lax pairs
is raised as a further question.  Class 6 is also identified with the
eleven-vertex model in Ref.~\cite{CorcoranDeLeeuw2024}; its associated
Landau--Lifshitz hierarchy was studied in
Ref.~\cite{LevinOlshanetskyZotov2014}.  These results provide the setting
for examining how the continuum time component follows from the
quantum lattice Lax equation.

The coherent-state lower symbol of an operator is its matrix element
in a product of spin coherent states, with any auxiliary-space indices
retained.  The spatial continuum Lax matrix follows from the leading
coefficient of a normalized lattice operator.  For the time component,
the lattice zero-curvature equation contains products whose factors
act on a common physical site.  Their lower symbols include a
contribution beyond the product of the separate symbols.  This
nonmultiplicativity is part of coherent-state symbol calculus
\cite{Perelomov1986,Schlichenmaier2010}.  In the continuum limit considered
here, the local spin representation remains fixed while the lattice
spacing tends to zero.  It is therefore necessary to determine which
operator-product contributions survive at the order of the continuum
time evolution.

In this paper we evaluate these contributions directly for both chains
and absorb them into a local correction of the continuum time Lax
matrix.  The Class 5 contribution is a total spatial derivative.  The
Class 6 calculation also retains a commutator involving the second
coefficient of the spatial lattice operator.  We derive the Hamiltonian
equations independently and give off-shell identities, with explicit
inverses, relating all their components to the Lax curvature.  For
Class 5, a separate classical $r$-matrix and monodromy construction
provides a comparison with the time matrix obtained from the lattice.

Section~\ref{sec:setup} fixes the lattice conventions and formulates the
continuum problem.  Sections~\ref{sec:class5-model} and
\ref{sec:class6-model} carry out the two model calculations.
Section~\ref{sec:comparison} compares the surviving corrections and
their dependence on the lattice operators.  We summarize the results
and discuss their scope and possible extensions in
Section~\ref{sec:summary-discussion}.
